\chapter{Monte Carlo Methods and Sign Problems}
\label{ch:constructive-monte-carlo-sign-problems}

\ConventionsEuclidean

Monte Carlo methods evaluate finite-regulator Euclidean expectation values by
sampling a probability distribution.  Their mathematical content is a
finite-dimensional integration problem plus a convergence theorem for a
Markov chain.  Their limitation in many QFTs is the absence of a positive
sampling weight compatible with the target observable.

\section{Finite-Regulator Expectation Values}

Let \(\Omega_\Lambda\) be a finite-dimensional regulator space and let
\[
  d\nu_\Lambda(\Phi)=\frac{1}{Z_\Lambda}e^{-S_\Lambda(\Phi)}d\Phi
\]
be a probability measure on it.  For an observable \(O:\Omega_\Lambda\to\mathbb C\),
\[
  \langle O\rangle_\Lambda
  =
  \int_{\Omega_\Lambda}O(\Phi)\,d\nu_\Lambda(\Phi).
\]
On a compact gauge lattice, \(d\Phi\) is a product of Haar measures and
finite-dimensional fermionic integrals have already been performed into a
determinant or Pfaffian when a purely bosonic sampling problem is used.

\section{Markov Chains and Detailed Balance}

A Markov transition kernel \(P(\Phi,d\Phi')\) preserves \(d\nu_\Lambda\) if
\[
  \int_{\Omega_\Lambda}d\nu_\Lambda(\Phi)\,P(\Phi,A)
  =
  \nu_\Lambda(A)
\]
for every measurable set \(A\).  Detailed balance is the stronger condition
\[
  d\nu_\Lambda(\Phi)P(\Phi,d\Phi')
  =
  d\nu_\Lambda(\Phi')P(\Phi',d\Phi).
\]
If the chain is irreducible and aperiodic in the finite-state case, or
satisfies the corresponding Harris recurrence hypotheses in a continuous
state space, the estimator
\[
  \overline O_N=\frac1N\sum_{n=1}^N O(\Phi_n)
\]
converges almost surely to \(\langle O\rangle_\Lambda\).

\section{Metropolis Updates as a Finite-Regulator Theorem}

The simplest exact check of the preceding definitions is a finite spin system.
Let \(\Lambda_L=(\mathbb Z/L\mathbb Z)^2\) be the periodic square lattice and
let a configuration be a map
\[
  s:\Lambda_L\to\{-1,+1\}.
\]
For real \(J,h\), define
\[
  H_L(s)
  =
  -J\sum_{\langle xy\rangle}s_xs_y-h\sum_xs_x ,
  \qquad
  \pi_\beta(s)=Z_{\beta,L}^{-1}e^{-\beta H_L(s)} .
\]
Here \(\langle xy\rangle\) denotes an unoriented nearest-neighbor bond and
\(Z_{\beta,L}\) is the finite sum over all \(2^{L^2}\) configurations.

The single-spin Metropolis chain chooses a site \(x\) uniformly, proposes
\[
  s\longmapsto s^x,\qquad
  s^x_y=
  \begin{cases}
    -s_x,& y=x,\\
    s_y,& y\neq x,
  \end{cases}
\]
and accepts the proposal with probability
\[
  a(s,x)=\min\{1,e^{-\beta(H_L(s^x)-H_L(s))}\}.
\]
The energy difference is local:
\[
  H_L(s^x)-H_L(s)
  =
  2s_x\left(J\sum_{y:\langle xy\rangle}s_y+h\right).
\]

\begin{proposition}[Detailed balance for the finite Ising Metropolis chain]
\label{prop:ising-metropolis-detailed-balance}
The transition kernel \(P\) defined by the preceding update satisfies
\[
  \pi_\beta(s)P(s,s')=\pi_\beta(s')P(s',s)
\]
for all configurations \(s,s'\).  If \(0<\beta<\infty\), the chain is
irreducible.  If additionally \((J,h)\neq(0,0)\), it is aperiodic on the
finite set \(\{-1,+1\}^{\Lambda_L}\).  For the degenerate infinite-temperature
Hamiltonian \(J=h=0\), adding a nonzero holding probability gives the usual
lazy aperiodic variant without changing the stationary measure.
\end{proposition}

\begin{proof}
If \(s'\) is not equal to \(s\) or to a one-spin flip of \(s\), both sides of
the detailed-balance identity vanish.  For \(s'=s^x\), write
\(\Delta H=H_L(s^x)-H_L(s)\).  The off-diagonal transition probabilities are
\[
  P(s,s^x)=\frac1{L^2}\min\{1,e^{-\beta\Delta H}\},
  \qquad
  P(s^x,s)=\frac1{L^2}\min\{1,e^{\beta\Delta H}\}.
\]
If \(\Delta H\geq0\), then
\[
  e^{-\beta H_L(s)}P(s,s^x)
  =
  \frac1{L^2}e^{-\beta H_L(s)}e^{-\beta\Delta H}
  =
  \frac1{L^2}e^{-\beta H_L(s^x)}
  =
  e^{-\beta H_L(s^x)}P(s^x,s).
\]
The case \(\Delta H<0\) is the same identity with \(s\) and \(s^x\)
interchanged.  Dividing by the finite partition function gives detailed
balance.

Irreducibility follows because any configuration can be reached from any
other by flipping the finite set of sites on which they differ, and every
single spin flip has positive proposal probability and positive acceptance
probability.  If \((J,h)\neq(0,0)\), there is a configuration \(s\) and a site
\(x\) with \(H_L(s^x)>H_L(s)\), so the chain has a positive holding
probability at \(s\).  In a finite irreducible Markov chain, one state with a
self-loop forces every state to have period one.  When \(J=h=0\), every
proposal is accepted and the one-spin-flip walk on the hypercube is bipartite;
the standard lazy modification inserts a holding probability and restores
aperiodicity.
\end{proof}

The companion script
\[
  \texttt{qft\_scripts/ising2d\_metropolis.py}
\]
implements exactly this finite chain and reports the acceptance rate, energy
per site, absolute magnetization, and a windowed estimate of
\(\tau_{\mathrm{int}}\).  Its smoke mode is a finite-regulator check of the
algorithmic identities above.  It does not estimate critical exponents, prove
the Onsager scaling limit, or construct the continuum Ising CFT.
The deterministic calculation check
\[
  \texttt{calculation-checks/ising\_metropolis\_finite\_checks.py}
\]
enumerates the \(2\times2\) periodic chain and verifies both the local energy
difference and the detailed-balance identity used by the script.

\section{Single-Link Metropolis for Finite Gauge Fields}

The same finite-regulator theorem should be stated for gauge variables,
because lattice gauge Monte Carlo samples compact group data rather than
coordinate fields.  For the \(\mathbb Z_2\) gauge model of
Proposition~\ref{prop:z2-gauge-character-expansion}, write
\[
  Q(U)=\sum_{p}U_p,\qquad
  \pi_\beta(U)=Z_\Lambda(\beta)^{-1}e^{\beta Q(U)} .
\]
For a link \(e\), let \(U^e\) be the configuration obtained by flipping
\(U_e\).  The only plaquettes whose values change are the plaquettes
containing \(e\).  If
\[
  B_e(U)=\sum_{p\supset e} U_p ,
\]
then
\[
  Q(U^e)-Q(U)=-2B_e(U).
\]
The single-link Metropolis chain chooses a link uniformly and accepts
\(U\mapsto U^e\) with probability
\[
  a(U,e)=\min\{1,\exp(\beta[Q(U^e)-Q(U)])\}.
\]

\begin{proposition}[Detailed balance for the finite \(\mathbb Z_2\) gauge chain]
\label{prop:z2-gauge-metropolis-detailed-balance}
The preceding single-link update satisfies
\[
  \pi_\beta(U)P(U,U')
  =
  \pi_\beta(U')P(U',U)
\]
for all finite lattice gauge configurations \(U,U'\).  On a connected finite
lattice, the chain is irreducible on the full link-configuration space.
Gauge-invariant observables may be averaged either on this full space or
after quotienting by gauge transformations; the full-space chain preserves
the gauge-invariant expectation values because \(\pi_\beta\) is gauge
invariant.
\end{proposition}

\begin{proof}
If \(U'\) differs from \(U\) by more than one link, the off-diagonal
transition probabilities in both directions vanish.  If \(U'=U^e\), set
\(\Delta Q=Q(U^e)-Q(U)\).  The off-diagonal transition probabilities are
\[
  P(U,U^e)=\frac1{|E|}\min\{1,e^{\beta\Delta Q}\},
  \qquad
  P(U^e,U)=\frac1{|E|}\min\{1,e^{-\beta\Delta Q}\}.
\]
For \(\Delta Q\geq0\),
\[
  e^{\beta Q(U)}P(U,U^e)
  =
  \frac1{|E|}e^{\beta Q(U)}
  =
  \frac1{|E|}e^{\beta Q(U^e)}e^{-\beta\Delta Q}
  =
  e^{\beta Q(U^e)}P(U^e,U).
\]
The case \(\Delta Q<0\) is the same identity with \(U\) and \(U^e\)
interchanged.  Division by \(Z_\Lambda(\beta)\) gives detailed balance.

Irreducibility is immediate on the finite link hypercube: any configuration
can be reached from any other by flipping the finite set of links on which
they differ, and each proposed flip has positive acceptance probability for
finite \(\beta\).  Gauge invariance follows because each \(U_p\) and each
closed Wilson loop is unchanged under \(U_{xy}\mapsto g_xU_{xy}g_y\).
Therefore averaging over the full link space includes gauge orbits with equal
weights and gives the same gauge-invariant expectations as any exact
gauge-orbit average with the corresponding orbit weights.
\end{proof}

The companion script
\[
  \texttt{qft\_scripts/z2\_gauge\_3d\_metropolis.py}
\]
implements this update for a periodic three-dimensional cubic lattice.  The
associated deterministic check
\[
  \texttt{calculation-checks/z2\_gauge\_metropolis\_checks.py}
\]
verifies the local score change, pairwise detailed balance, gauge invariance
of the plaquette action and Wilson loops, and the equality between the
\(1\times1\) Wilson-loop average and the plaquette average.

\section{Autocorrelation and Effective Sample Size}

For a stationary chain, define the connected autocorrelation function
\[
  C_O(t)
  =
  \langle O(\Phi_t)O(\Phi_0)\rangle
  -\langle O\rangle^2 .
\]
The integrated autocorrelation time is
\[
  \tau_{\mathrm{int}}(O)
  =
  \frac12+\sum_{t=1}^\infty\frac{C_O(t)}{C_O(0)}
\]
when the sum converges.  The variance of the sample mean is then
\[
  \operatorname{Var}(\overline O_N)
  =
  \frac{2\tau_{\mathrm{int}}(O)}{N}
  \operatorname{Var}(O)
  +o(N^{-1}).
\]
Critical slowing down is the growth of \(\tau_{\mathrm{int}}\) along a scaling
trajectory toward a continuum limit.

\section{Reweighting and the Sign Problem}

Suppose the target finite-regulator integral has complex weight
\[
  e^{-S(\Phi)}=e^{-S_R(\Phi)}e^{\ii\Theta(\Phi)}
\]
with positive reference measure
\[
  d\nu_R=\frac{1}{Z_R}e^{-S_R(\Phi)}d\Phi .
\]
Then
\[
  \langle O\rangle
  =
  \frac{\langle Oe^{\ii\Theta}\rangle_R}
  {\langle e^{\ii\Theta}\rangle_R}.
\]
The denominator is the average phase.  If the free-energy densities of the
target and phase-quenched systems differ by \(\Delta f\) in volume \(V\), then
\[
  \left|\langle e^{\ii\Theta}\rangle_R\right|
  =
  \exp[-V\Delta f].
\]
Estimating a ratio with a denominator of this size requires a number of
independent samples growing as \(\exp[2V\Delta f]\) for fixed relative error,
unless the numerator and denominator are computed by a method that avoids
phase reweighting.

\section{Fermion Determinants}

For lattice fermions with quadratic action
\[
  \bar\psi D[U]\psi ,
\]
Berezin integration gives
\[
  \int d\bar\psi\,d\psi\,e^{-\bar\psi D[U]\psi}
  =
  \det D[U].
\]
If \(D\) obeys a positivity relation, such as
\[
  \gamma_5D\gamma_5=D^\dagger
\]
together with flavor pairing, the determinant can be nonnegative.  At finite
chemical potential or in chiral gauge theories this positivity may fail.  The
Monte Carlo problem is then not merely inefficient sampling; the finite
regulator has no declared positive probability measure whose expectations are
the target correlators without reweighting or contour deformation.

\section{Continuum Interpretation}

Monte Carlo data become QFT data only after the regulator, source
normalization, operator renormalization, scaling trajectory, and error limits
are specified.  A numerical estimate at finite lattice spacing is a statement
about \(d\nu_\Lambda\).  A continuum claim requires convergence of a family of
finite-regulator expectations in the topology needed for OS reconstruction,
Wightman reconstruction, or the intended observable sector.

\begin{openproblem}[Sign-problem-preserving continuum construction]
\label{op:sign-problem-continuum-construction}
Classify QFTs with complex Euclidean weights for which a nonperturbative
continuum limit can be constructed by a positive enlarged-variable measure,
an exact contour deformation, or another finite-regulator method with
controlled variance and locality.
\end{openproblem}
